% GENERATED FILE. Edit canonical lessons, then run npm run generate:books.
% canonical-source-hash: fnv1a64:93819ef3efde229e
% canonical-lessons: TA-C333-porkkalam, TA-C333-porvirar, TA-C333-val-sword, TA-C333-ketayam, TA-C333-kavacam

\chapter{The Battlefield}
\label{ch:ta-the-battlefield}
\hlchaptermodality{333}

\begin{chapteropening}
\textbf{By the end of this chapter:} \emph{I can describe soldiers and what they carried.}

Complete the last lesson of chapter 333: I can describe soldiers and what they carried.
\end{chapteropening}

\section[\texorpdfstring{\ta{போர்க்களம்} (pōrkkaḷam)}{போர்க்களம் (pōrkkaḷam)}]{\ta{போர்க்களம்} (pōrkkaḷam) — a battlefield}

\label{lesson:TA-C333-porkkalam}



\begin{quote}
Before the new one: say the Tamil for to attack, then the Tamil for a navy.
\end{quote}

\subsection*{You'll want to know: \ta{போர்க்களம்}}
\textbf{\ta{போர்க்களம்}} — \emph{pōrkkaḷam} — \textquotedblleft{}a battlefield\textquotedblright{}.

\begin{quote}
\ta{போர்க்களத்தில்} \ta{பல} \ta{வீரர்கள்} \ta{இறந்தார்கள்}.
\end{quote}

\begin{quote}
\emph{pōrkkaḷattil pala vīrarkaḷ iṟantārkaḷ.}
\end{quote}

\begin{quote}
Many soldiers died on the battlefield.
\end{quote}

\ta{களம்}, field, open ground.

\begin{cousinweb}
\textbf{\ta{போர்க்களம்}} (\emph{pōrkkaḷam}) is a Tamil compound built from \ta{போர்}, \textquotedblleft{}battle\textquotedblright{} (DEDR 4540), and \ta{களம்}, \textquotedblleft{}open space, battlefield\textquotedblright{} (DEDR 1376).
\end{cousinweb}

\begin{grammarlens}[title={What you've built}]
One more word for telling what happened.
\end{grammarlens}

\subsection*{Your turn}
\begin{itemize}
\raggedright
  \item \emph{Say it:} \emph{pōrkkaḷam}
  \item \emph{Say it:} \emph{pōrkkaḷam}, once more
  \item \emph{Say it:} \emph{pōrkkaḷattil pala vīrarkaḷ iṟantārkaḷ.} — Many soldiers died on the battlefield.
  \item \emph{Recall:} say the Tamil for to attack, then the Tamil for a navy, then say \emph{pōrkkaḷam} again
\end{itemize}

\begin{tcolorbox}[breakable,colback=teal!4,colframe=teal!35!black,title={Before you move on}]
What does \ta{போர்க்களம்} mean? (\textquotedblleft{}A battlefield\textquotedblright{}.) Say it once more.
\end{tcolorbox}
\section[\texorpdfstring{\ta{போர்வீரர்} (pōrvīrar)}{போர்வீரர் (pōrvīrar)}]{\ta{போர்வீரர்} (pōrvīrar) — a soldier; a warrior}

\label{lesson:TA-C333-porvirar}



\begin{quote}
Before the new one: say the Tamil for a navy, then the Tamil for a battlefield.
\end{quote}

\subsection*{You'll want to know: \ta{போர்வீரர்}}
\textbf{\ta{போர்வீரர்}} — \emph{pōrvīrar} — \textquotedblleft{}a soldier; a warrior\textquotedblright{}.

\begin{quote}
\ta{போர்வீரர்கள்} \ta{வரிசையாக} \ta{நின்றார்கள்}.
\end{quote}

\begin{quote}
\emph{pōrvīrarkaḷ varicaiyāka niṉṟārkaḷ.}
\end{quote}

\begin{quote}
The soldiers stood in rows.
\end{quote}

Plural \ta{போர்வீரர்கள்}.

\begin{cousinweb}
\textbf{\ta{போர்வீரர்}} (\emph{pōrvīrar}) joins native \ta{போர்}, \textquotedblleft{}battle\textquotedblright{} (DEDR 4540), with \ta{வீரர்}, from Sanskrit \emph{vīra}, \textquotedblleft{}man, hero\textquotedblright{} (CDIAL 12056).
\end{cousinweb}

\begin{grammarlens}[title={What you've built}]
One more word for telling what happened.
\end{grammarlens}

\subsection*{Your turn}
\begin{itemize}
\raggedright
  \item \emph{Say it:} \emph{pōrvīrar}
  \item \emph{Say it:} \emph{pōrvīrar}, once more
  \item \emph{Say it:} \emph{pōrvīrarkaḷ varicaiyāka niṉṟārkaḷ.} — The soldiers stood in rows.
  \item \emph{Recall:} say the Tamil for a navy, then the Tamil for a battlefield, then say \emph{pōrvīrar} again
\end{itemize}

\begin{tcolorbox}[breakable,colback=teal!4,colframe=teal!35!black,title={Before you move on}]
What does \ta{போர்வீரர்} mean? (\textquotedblleft{}A soldier; a warrior\textquotedblright{}.) Say it once more.
\end{tcolorbox}
\section[\texorpdfstring{\ta{வாள்} (vāḷ)}{வாள் (vāḷ)}]{\ta{வாள்} (vāḷ) — a sword}

\label{lesson:TA-C333-val-sword}



\begin{quote}
Before the new one: say the Tamil for a battlefield, then the Tamil for a soldier.
\end{quote}

\subsection*{You'll want to know: \ta{வாள்}}
\textbf{\ta{வாள்}} — \emph{vāḷ} — \textquotedblleft{}a sword\textquotedblright{}.

\begin{quote}
\ta{மன்னர்} \ta{தன்} \ta{வாளை} \ta{உயர்த்தினார்}.
\end{quote}

\begin{quote}
\emph{maṉṉar taṉ vāḷai uyarttiṉār.}
\end{quote}

\begin{quote}
The king raised his sword.
\end{quote}

Plural \ta{வாள்கள்}.

\begin{cousinweb}
\textbf{\ta{வாள்}} (\emph{vāḷ}) is native Dravidian (DEDR 5376): sword, saw, sharpness.
\end{cousinweb}

\begin{grammarlens}[title={What you've built}]
One more word for telling what happened.
\end{grammarlens}

\subsection*{Your turn}
\begin{itemize}
\raggedright
  \item \emph{Say it:} \emph{vāḷ}
  \item \emph{Say it:} \emph{vāḷ}, once more
  \item \emph{Say it:} \emph{maṉṉar taṉ vāḷai uyarttiṉār.} — The king raised his sword.
  \item \emph{Recall:} say the Tamil for a battlefield, then the Tamil for a soldier, then say \emph{vāḷ} again
\end{itemize}

\begin{tcolorbox}[breakable,colback=teal!4,colframe=teal!35!black,title={Before you move on}]
What does \ta{வாள்} mean? (\textquotedblleft{}A sword\textquotedblright{}.) Say it once more.
\end{tcolorbox}
\section[\texorpdfstring{\ta{கேடயம்} (kēṭayam)}{கேடயம் (kēṭayam)}]{\ta{கேடயம்} (kēṭayam) — a shield}

\label{lesson:TA-C333-ketayam}



\begin{quote}
Before the new one: say the Tamil for a soldier, then the Tamil for a sword.
\end{quote}

\subsection*{You'll want to know: \ta{கேடயம்}}
\textbf{\ta{கேடயம்}} — \emph{kēṭayam} — \textquotedblleft{}a shield\textquotedblright{}.

\begin{quote}
\ta{அவன்} \ta{கேடயத்தால்} \ta{தன்னைக்} \ta{காத்துக்கொண்டான்}.
\end{quote}

\begin{quote}
\emph{avaṉ kēṭayattāl taṉṉaik kāttukkoṇṭāṉ.}
\end{quote}

\begin{quote}
He protected himself with a shield.
\end{quote}

-\ta{ஆல்} marks the means: \ta{கேடயத்தால்}, with a shield.

\begin{cousinweb}
\textbf{\ta{கேடயம்}} (\emph{kēṭayam}) is borrowed from Sanskrit \emph{kheṭaka}, \textquotedblleft{}a shield\textquotedblright{} (Monier-Williams).
\end{cousinweb}

\begin{grammarlens}[title={What you've built}]
One more word for telling what happened.
\end{grammarlens}

\subsection*{Your turn}
\begin{itemize}
\raggedright
  \item \emph{Say it:} \emph{kēṭayam}
  \item \emph{Say it:} \emph{kēṭayam}, once more
  \item \emph{Say it:} \emph{avaṉ kēṭayattāl taṉṉaik kāttukkoṇṭāṉ.} — He protected himself with a shield.
  \item \emph{Recall:} say the Tamil for a soldier, then the Tamil for a sword, then say \emph{kēṭayam} again
\end{itemize}

\begin{tcolorbox}[breakable,colback=teal!4,colframe=teal!35!black,title={Before you move on}]
What does \ta{கேடயம்} mean? (\textquotedblleft{}A shield\textquotedblright{}.) Say it once more.
\end{tcolorbox}
\section[\texorpdfstring{\ta{கவசம்} (kavacam)}{கவசம் (kavacam)}]{\ta{கவசம்} (kavacam) — armour}

\label{lesson:TA-C333-kavacam}



\begin{quote}
Before the new one: say the Tamil for a sword, then the Tamil for a shield.
\end{quote}

\subsection*{You'll want to know: \ta{கவசம்}}
\textbf{\ta{கவசம்}} — \emph{kavacam} — \textquotedblleft{}armour\textquotedblright{}.

\begin{quote}
\ta{வீரர்கள்} \ta{இரும்புக்} \ta{கவசம்} \ta{அணிந்திருந்தார்கள்}.
\end{quote}

\begin{quote}
\emph{vīrarkaḷ irumpuk kavacam aṇintiruntārkaḷ.}
\end{quote}

\begin{quote}
The warriors wore iron armour.
\end{quote}

Also any protective cover: \ta{தலைக்கவசம்}, a helmet.

\begin{cousinweb}
\textbf{\ta{கவசம்}} (\emph{kavacam}) is borrowed from Sanskrit \emph{kavaca}, \textquotedblleft{}a coat of mail\textquotedblright{} (CDIAL 2957); CDIAL notes a possible Dravidian source for the Sanskrit word itself.
\end{cousinweb}

\begin{grammarlens}[title={What you've built}]
One more word for telling what happened.
\end{grammarlens}

\subsection*{Your turn}
\begin{itemize}
\raggedright
  \item \emph{Say it:} \emph{kavacam}
  \item \emph{Say it:} \emph{kavacam}, once more
  \item \emph{Say it:} \emph{vīrarkaḷ irumpuk kavacam aṇintiruntārkaḷ.} — The warriors wore iron armour.
  \item \emph{Recall:} say the Tamil for a sword, then the Tamil for a shield, then say \emph{kavacam} again
\end{itemize}

\begin{tcolorbox}[breakable,colback=teal!4,colframe=teal!35!black,title={Before you move on}]
What does \ta{கவசம்} mean? (\textquotedblleft{}Armour\textquotedblright{}.) Say it once more.
\end{tcolorbox}
